\documentclass[11pt]{article}

\usepackage[utf8]{inputenc}
\usepackage[T1]{fontenc}
\usepackage{lmodern}
\usepackage{amsmath,amssymb,amsfonts,amsthm}
\usepackage{bm}
\usepackage{mathtools}
\usepackage{geometry}
\usepackage{booktabs}
\usepackage{hyperref}
\usepackage{microtype}

\geometry{margin=1in}

\newtheorem{example}{Example}
\newtheorem{assumption}{Assumption}
\newtheorem{remark}{Remark}
\newtheorem{lemma}{Lemma}

\title{Theoretical Basis of the SPSA Process}
\author{}
\date{}

\begin{document}
\maketitle

\section{The SPSA ``process''}

We analyze an elegant chess–oriented version of the SPSA algorithm invented by
Stockfish author Joona Kiiski~\cite{Kiiski}.

Let $(\bm{c}_i)_{i \ge 0} \subset \mathbb{R}^n$ be a sequence of column vectors and let $(r_i)_{i \ge 0} \subset \mathbb{R}_{>0}$. To
go from a vector of engine parameters


\[
  \bm{\theta}_i = (\theta_i^{(j)})_{j=1}^n \in \mathbb{R}^n
\]


to a new vector $\bm{\theta}_{i+1} \in \mathbb{R}^n$ we play a match of $\bm{\theta}_i + \bm{c}_i$ against $\bm{\theta}_i - \bm{c}_i$ and put
\begin{equation}
  \bm{\theta}_{i+1} =
  \begin{cases}
    \bm{\theta}_i + r_i \bm{c}_i & \text{in case of a win},\

\[4pt]
    \bm{\theta}_i                & \text{in case of a draw},\

\[4pt]
    \bm{\theta}_i - r_i \bm{c}_i & \text{in case of a loss.}
  \end{cases}
  \label{eq:update}
\end{equation}

Since we have biased\footnote{For the purpose of exposition we assume that the games are played with an opening book that
is perfectly balanced. In practice this will never be satisfied but one may simulate it by replaying
every game with reversed colors. This will be assumed silently below.}
the update of $\bm{\theta}_i$ in the direction of the engine that won the
game, the hope is that $\bm{\theta}_i$ for $i > 0$ will yield a stronger engine than $\bm{\theta}_0$.

The next two examples illustrate standard choices for $\bm{c}_i$.

\begin{example}
Let $c_0,\dots,c_{n-1} > 0$ and put


\[
  \bm{c}_i = c_{i \bmod n} \,\bm{e}_{i \bmod n},
\]


where $\bm{e}_j$ is the $j$th standard basis vector of $\mathbb{R}^n$. This is a variant of the FDSA algorithm. We have


\[
  \mathbb{E}(\bm{c}_i \bm{c}_i^\top) = \operatorname{diag}(c_0^2,\dots,c_{n-1}^2).
\]


\end{example}

\begin{example}
Let $c_0,\dots,c_{n-1} > 0$ and put


\[
  \bm{c}_i = \sum_{j=0}^{n-1} \alpha_{i,j} c_j \bm{e}_j
\]


where each $\alpha_{i,j}$ is chosen uniformly at random in $\{\pm 1\}$. Then


\[
  \mathbb{E}(\bm{c}_i \bm{c}_i^\top) = \operatorname{diag}(c_0^2,\dots,c_{n-1}^2).
\]


This is the SPSA algorithm.
\end{example}

\begin{assumption}
We allow $(\bm{c}_i)_{i \ge 0}$ to be stochastic but throughout we assume that
$(\bm{c}_i \bm{c}_i^\top)_i$ stays close to its short–term running average / expectation value (denoted by
$\mathbb{E}(\bm{c}_i \bm{c}_i^\top)$) and that the latter is an invertible $n \times n$ matrix.
\end{assumption}

\section{Main results}

\subsection{Statements}

Let $e$ be the function which maps $\bm{\theta} \in \mathbb{R}^n$ to Elo. Assume
that $e$ is (approximately) quadratic and attains its maximum value at $\bm{\theta} = \bm{\theta}^\ast$. We
normalize $e$ such that $e(\bm{\theta}^\ast) = 0$.

Apply the generalized SPSA algorithm with constant $r$ and with $\bm{c}_i$ satisfying
Assumption~\ref{ass:Ecct} below such that $\mathbb{E}(\bm{c}_i \bm{c}_i^\top)$ is constant. We will heuristically derive some
reasonable approximations.

\begin{assumption}
\label{ass:Ecct}
For all $i$, $\mathbb{E}(\bm{c}_i \bm{c}_i^\top)$ exists, is constant in $i$, and is invertible.
\end{assumption}

\begin{enumerate}
\item Asymptotically the expectation value of $e(\bm{\theta}_i)$ will satisfy
\begin{equation}
  \mathbb{E}\bigl(e(\bm{\theta}_i)\bigr)
  \;\approx\;
  - \frac{r}{8}\, C\, n (1-d),
  \label{eq:main1}
\end{equation}
with $d \in [0,1]$ being equal to the draw ratio\footnote{With an unbalanced book and replayed color–reversed games the factor $1-d$ in the formula
can be reduced to $1-d-4b^2$ where $b$ is the RMS bias of the opening book (computed using
$0, 1/2, 1$ scoring, see~\cite{AccountingIdentity}). The same holds true for similar formulas below.}
and


\[
  C = \frac{800}{\ln 10} \approx 347.43558552.
\]



\item Let $0 \le \gamma \le 1$ and define $z$ by


\[
  \gamma = \mathbb{P}(X < z) \quad\text{for}\quad X \sim \chi^2_n.
\]


For example, for $\gamma = 0.95$ we have in low dimensions:
\begin{center}
\begin{tabular}{@{}cc@{}}
\toprule
\textbf{dim} & \textbf{$z$} \\
\midrule
1 & 3.8414588206941236 \\
2 & 5.9914645471079799 \\
3 & 7.8147279032511765 \\
4 & 9.4877290367811540 \\
\bottomrule
\end{tabular}
\end{center}

With probability $\gamma$, $e(\bm{\theta}_i)$ will asymptotically satisfy
\begin{equation}
  e(\bm{\theta}_i)
  \;\gtrsim\;
  - \frac{r}{8}\, C\, z (1-d).
  \label{eq:main2}
\end{equation}

\item The expectation value of $\bm{\theta}_i$ will converge linearly towards $\bm{\theta}^\ast$ as
\begin{equation}
  \|\mathbb{E}(\bm{\theta}_i - \bm{\theta}^\ast)\|
  \approx
  \|\bm{\theta}_0 - \bm{\theta}^\ast\| \, e^{-i/\lambda}
  \label{eq:main3}
\end{equation}
with the time constant $\lambda$ bounded by
\begin{equation}
  \lambda \;\approx\; \frac{C}{2 r \mu},
  \label{eq:lambda}
\end{equation}
where $\mu$ is the smallest eigenvalue of


\[
  \mathbb{E}(\bm{c}_i \bm{c}_i^\top)\, \operatorname{Hess}(e)(\bm{\theta}^\ast).
\]


\end{enumerate}

\begin{remark}
We see from (3) that the nicest possible case is when


\[
  \mathbb{E}(\bm{c}_i \bm{c}_i^\top)\, \operatorname{Hess}(e)(\bm{\theta}^\ast) = \mu I_{n}.
\]


It follows from the derivation below that in that case the update~\eqref{eq:update} is (on average)
proportional to


\[
  r\, \operatorname{Hess}(e)(\bm{\theta}^\ast)^{-1} \bigl(\nabla e\bigr)(\bm{\theta}_i),
\]


which is in particular the optimal direction for gradient descent.
\end{remark}

\subsection{Derivations}

Below vectors are always column vectors. We will also switch
to continuous time. Hence we write $\bm{\theta}(t)$ (or simply $\bm{\theta}$) instead of $\bm{\theta}_i$, etc.

We can think of the dynamics of $\bm{\theta}$ as a discrete version of the continuous process
\begin{equation}
  d\bm{\theta}_t
  = r \bm{c}\,\bigl(\mu_{\theta,\bm{c}}\,dt + \sigma_{\theta,\bm{c}}\,dW_t\bigr),
\end{equation}
where $(\mu_{\theta,\bm{c}}, \sigma_{\theta,\bm{c}})$ are the mean and standard deviation of a single game between
parameters $\bm{\theta} \pm \bm{c}$ (scored as $-1,0,1$) and $W_t$ is a Wiener process (Brownian motion).
We simplify further and assume that $\sigma_{\theta,\bm{c}}$ is independent of $\bm{\theta}, \bm{c}$ so that
\begin{equation}
  \sigma_{\theta,\bm{c}}^2 = 1 - d,
  \label{eq:sigma}
\end{equation}
where $d$ is the draw ratio.

Let $f(\bm{\theta})$ be the function that maps $\bm{\theta}$ to the expected score against a reference
engine (scored as $-1,0,1$). Assuming linearity we get
\begin{align*}
  \mu_{\theta,\bm{c}}
  &= \operatorname{score}(\bm{\theta}+\bm{c},\bm{\theta}-\bm{c}) \\
  &= f(\bm{\theta}+\bm{c}) - f(\bm{\theta}-\bm{c})
   = 2 \bm{c}^\top (\nabla f)(\bm{\theta}),
\end{align*}
and hence
\begin{equation}
  d\bm{\theta}_t
  = 2 r \bm{c}\bm{c}^\top (\nabla f)(\bm{\theta})\,dt
    + r \sigma_{\theta,\bm{c}} \bm{c}\,dW_t.
\end{equation}

We now replace $\bm{c}\bm{c}^\top$ by its short–term average / expectation value which we denote
by $\mathbb{E}(\bm{c}\bm{c}^\top)$ and we put


\[
  A = r\,\mathbb{E}(\bm{c}\bm{c}^\top).
\]


Note that $A$ is symmetric. We obtain
\begin{equation}
  d\bm{\theta}_t
  = 2 A (\nabla f)(\bm{\theta})\,dt + r \sigma_{\theta,\bm{c}} \bm{c}\,dW_t.
\end{equation}

The conversion between $f$ and $e$ is approximately given by
\begin{equation}
  f(\bm{\theta}) = \frac{e(\bm{\theta})}{C}.
  \label{eq:f-e}
\end{equation}
So $f$ is approximately quadratic as well. Hence (translating such that $\bm{\theta}^\ast = \bm{0}$)
we may write


\[
  f(\bm{\theta}) = -\tfrac{1}{2} \bm{\theta}^\top E \bm{\theta},
\]


with $E$ a positive–definite symmetric $n \times n$ matrix. So


\[
  (\nabla f)(\bm{\theta}) = - E \bm{\theta},
\]


from which we get
\begin{equation}
  d\bm{\theta}_t
  = -2 A E \bm{\theta}\,dt + r \sigma_{\theta,\bm{c}} \bm{c}\,dW_t.
\end{equation}

The solution of the corresponding homogeneous equation


\[
  d\bm{\theta}_t = -2 A E \bm{\theta}\,dt
\]


is given by


\[
  \bm{\theta}(t) = \exp(-2 A E t)\,\bm{D},
\]


where $\bm{D}$ is a constant column vector. We now solve the inhomogeneous system
using variation of constants (i.e.\ we make $\bm{D}$ depend on $t$). We get


\[
  \exp(-2 A E t)\, d\bm{D}(t) = r \sigma_{\theta,\bm{c}} \bm{c}\,dW_t,
\]


hence


\[
  \bm{D}(t)
  = \bm{D}(0)
    + r \sigma_{\theta,\bm{c}} \int_0^t \exp(2 A E s)\,\bm{c}\,dW_s.
\]


So
\begin{equation}
  \bm{\theta}(t)
  = \exp(-2 A E t)\,\bm{\theta}(0)
    + r \sigma_{\theta,\bm{c}} \int_0^t \exp(2 A E (s-t))\,\bm{c}\,dW_s.
\end{equation}

Thus $\bm{\theta}(t)$ is multivariate normal with expectation value


\[
  \mathbb{E}(\bm{\theta}(t))
  = \exp(-2 A E t)\,\bm{\theta}(0)
  = \exp\bigl(-2 r\,\mathbb{E}(\bm{c}\bm{c}^\top) E\,t\bigr)\,\bm{\theta}(0),
\]


from which we deduce the convergence behaviour in (3), using the conversion~\eqref{eq:f-e} combined with~\eqref{eq:sigma}.

We now calculate the covariance of $\bm{\theta}(t)$. Using the stochastic integral term and again replacing $\bm{c}\bm{c}^\top$ by its
average, we obtain
\begin{align*}
  \operatorname{Cov}(\bm{\theta}(t))
  &= r^2 \sigma_{\theta,\bm{c}}^2
     \int_0^t \exp(2 A E (s-t))\,\bm{c}\bm{c}^\top\,\exp(2 E A (s-t))\,ds \\
  &\approx r^2 \sigma_{\theta,\bm{c}}^2
     \int_0^t \exp(2 A E (s-t))\,\mathbb{E}(\bm{c}\bm{c}^\top)\,\exp(2 E A (s-t))\,ds \\
  &= r^2 \sigma_{\theta,\bm{c}}^2
     \int_0^t \exp\bigl(4 E A (s-t)\bigr)\,A\,ds.
\end{align*}

Under the simplifying assumption that $E$ and $A$ commute, we get


\[
  \operatorname{Cov}(\bm{\theta}(t))
  = \frac{r^2 \sigma_{\theta,\bm{c}}^2}{4}\,
    A^{-1} \bigl(I - \exp(-4 E A t)\bigr).
\]


Hence asymptotically


\[
  \operatorname{Cov}(\bm{\theta})
  = \frac{r^2 \sigma_{\theta,\bm{c}}^2}{4}\,A^{-1},
\]


and so asymptotically the distribution of $\bm{\theta}$ is


\[
  \bm{\theta} \sim \mathcal{N}\!\left(\bm{0},\,\frac{r^2 \sigma_{\theta,\bm{c}}^2}{4}\,A^{-1}\right).
\]



Recall the following.

\begin{lemma}
Let $E$ be a positive–definite symmetric $n \times n$ matrix and


\[
  \bm{X} \sim \mathcal{N}(\bm{0}, E).
\]


Then


\[
  \bm{X}^\top E^{-1} \bm{X} \sim \chi^2_n.
\]


\end{lemma}

\begin{proof}
We may write $E = P P^\top$ with $P$ invertible. Put $\bm{Y} = P^{-1} \bm{X}$. Then


\[
  \operatorname{Cov}(\bm{Y})
  = P^{-1} \operatorname{Cov}(\bm{X}) (P^{-1})^\top
  = P^{-1} E (P^{-1})^\top
  = P^{-1} P P^\top (P^{-1})^\top
  = I.
\]


Hence $\bm{Y} \sim \mathcal{N}(\bm{0}, I)$ and thus $\bm{Y}^\top \bm{Y} \sim \chi^2_n$. But


\[
  \bm{Y}^\top \bm{Y}
  = \bm{X}^\top (P^{-1})^\top P^{-1} \bm{X}
  = \bm{X}^\top E^{-1} \bm{X}.
\]


\end{proof}

From this lemma and the asymptotic covariance we obtain that, up to scaling, $e(\bm{\theta})$ follows (asymptotically) a chi–square law, which yields the qualitative behaviour stated in (1) and (2), once we translate back via~\eqref{eq:f-e} and~\eqref{eq:sigma}.

\subsection{The case of diagonal Hessian}

In this section we assume


\[
  e(\bm{\theta}) = -\sum_{j=1}^n \varepsilon_j \theta_j^2
\]


and furthermore that we are in the situation of Example~1.2 so that


\[
  \mathbb{E}(\bm{c}\bm{c}^\top) = \operatorname{diag}(c_0^2,\dots,c_{n-1}^2).
\]



Then Statement (3) may be refined as follows.

\begin{enumerate}
\setcounter{enumi}{3}
\item The expectation value of $\theta_j(t)$, $j = 1,\dots,n$, will converge linearly towards
$\theta_j^\ast$ as


\[
  \mathbb{E}\bigl(\theta_j(t)\bigr)
  = \theta_j^\ast
    + \bigl(\theta_j(0) - \theta_j^\ast\bigr) \exp\!\Bigl(-\frac{t}{\lambda_j}\Bigr),
\]


with the time constant $\lambda_j$ given by
\begin{equation}
  \lambda_j
  = \frac{C}{8 r c_j^2 \varepsilon_j}.
  \label{eq:lambda_j}
\end{equation}
\end{enumerate}

We will now discuss some more detailed results. We have that each $\theta_j(t)$ is normally
distributed, and moreover


\[
  \operatorname{Var}\bigl(\theta_j(t)\bigr)
  = \frac{r(1-d)C}{8\varepsilon_j}
    \Bigl(1 - \exp\bigl(-8 r c_j^2 \varepsilon_j t / C\bigr)\Bigr).
\]



In particular we obtain


\[
  \mathbb{E}\bigl(\theta_j(t)\bigr)
  = \exp\bigl(-4 r c_j^2 \varepsilon_j t / C\bigr)\,\theta_j(0),
\]


and


\[
  \mathbb{E}\bigl(e(\bm{\theta}(t))\bigr)
  = \sum_{j=1}^n \left(
      -\varepsilon_j \exp\bigl(-8 r c_j^2 \varepsilon_j t / C\bigr)\,\theta_j(0)^2
      + \frac{r(1-d)C}{8}\Bigl(1 - \exp\bigl(-8 r c_j^2 \varepsilon_j t / C\bigr)\Bigr)
    \right).
\]



Furthermore $e(\bm{\theta}(t))$ follows a generalized chi–square distribution~\cite{wiki:generalized-chi2}. The SPSA simula\-tor~\cite{spsa_simul}
contains routines to calculate the cumulative distribution function (CDF) and
the percent point function (PPF) for the generalized chi–square distribution. The CDF
routine was originally based on MATLAB code written by Abhranil Das~\cite{Das}
and is based on Ruben's algorithm~\cite{Ruben}. Essentially, in Ruben's algorithm one develops
the characteristic function $\varphi(t)$ of a generalized chi–square distribution as a power series in
$(1-2 i t \rho)^{-1/2}$ for a suitably chosen $\rho$. This leads to a description of the generalized
chi–square distribution as a mixture of ordinary scaled $\chi^2_{n-2i}$–distributions for $i \ge 0$.

It follows from Ruben's formulas that if


\[
  \exp\bigl(-8 r c_j^2 \varepsilon_j t / C\bigr)
  = \exp\bigl(-2t/\lambda_j\bigr)
\]


is small, then the distribution of $e(\bm{\theta}(t))$ satisfies approximately


\[
  \frac{n\,e(\bm{\theta}(t))}{\mathbb{E}(e(\bm{\theta}(t)))} \sim \chi^2_n.
\]



We have also shown that asymptotically


\[
  \frac{e(\bm{\theta})}{\mathbb{E}(e(\bm{\theta}))} \sim \chi^2_n,
\]


so


\[
  \frac{n\,e(\bm{\theta}(0))}{\mathbb{E}(e(\bm{\theta}(0)))}
  \quad\text{and}\quad
  \frac{e(\bm{\theta}(t))}{\mathbb{E}(e(\bm{\theta}(t)))}
\]


serve as scale factors to translate asymptotic location quantities (e.g.\ mean and
percentiles) into the corresponding quantities at time $t$ for $t \gg 0$.

Let us finally note that if $c_j^2 \varepsilon_j$ is independent of $j$ so that we can put $\lambda = \lambda_j$, then
we get the following elegant formula
\begin{equation}
  \mathbb{E}\bigl(e(\bm{\theta}(t))\bigr)
  = \mathbb{E}\bigl(e(\bm{\theta}(0))\bigr)\exp(-2t/\lambda)
    + \mathbb{E}\bigl(e(\bm{\theta})\bigr)\bigl(1-\exp(-2t/\lambda)\bigr).
  \label{eq:e-evolution}
\end{equation}

\subsection{Optimal approach}

Fix a target $-p < 0$. Starting from~\eqref{eq:e-evolution} we attempt to choose $r$ in
such a way that $\mathbb{E}\bigl(e(\bm{\theta}(t))\bigr) = -p$ is achieved for the smallest possible~$t$. To obtain
$t$ from $r$ we have to solve the following equation:


\[
  p = p_0 \exp(-a r t) + b r \bigl(1-\exp(-a r t)\bigr),
\]


with
\begin{equation}
  p_0 = -\mathbb{E}\bigl(e(\bm{\theta}(0))\bigr), \qquad
  a = \frac{8 c_j^2 \varepsilon_j}{C}, \qquad
  b = \frac{n(1-d) C}{8}.
  \label{eq:abc}
\end{equation}

After some rescaling to eliminate parameters we obtain the universal equation
\begin{equation}
  \exp(-R T) + R\bigl(1-\exp(-R T)\bigr) = p,
\end{equation}
with


\[
  R = \frac{b r}{p_0}, \qquad
  p = \frac{p}{p_0}, \qquad
  T = a t.
\]



Using~\eqref{eq:abc} we may express $T$ in terms of $R$ and $p$. Then we may determine $R$
such that $\partial T(R,p)/\partial R = 0$. It seems this last step must be done numerically. The
optimal $R$ as a function of $p$ has a characteristic shape: starting near $0$ at $p=0$, growing to a maximum,
and then decreasing towards $0$ as $p \to 1$.

\section*{References}

\begin{thebibliography}{9}

\bibitem{wiki:generalized-chi2}
Wikipedia contributors, \emph{Generalized chi-squared distribution},
\url{https://en.wikipedia.org/w/index.php?title=Generalized_chi-squared_distribution&oldid=958566492}.

\bibitem{Das}
A.~Das,
\emph{MATLAB toolbox for classifying amongst normal distributions},
\url{https://github.com/abhranildas/classify}.

\bibitem{Kiiski}
J.~Kiiski,
\emph{SPSA Tuner for Stockfish Chess Engine},
\url{https://github.com/zamar/spsa}.

\bibitem{Ruben}
H.~Ruben,
Probability content of regions under spherical normal distributions. IV. The distribution of homogeneous and non-homogeneous quadratic functions of normal variables,
\emph{Ann. Math. Statist.} 33 (1962), 542--570.

\bibitem{spsa_simul}
M.~Van den Bergh,
\emph{A multi threaded simulator for the chess version of the SPSA algorithm
created by Joona Kiiski},
\url{https://github.com/vdbergh/spsa_simul}.

\bibitem{AccountingIdentity}
\emph{The accounting identity},
\url{http://hardy.uhasselt.be/Fishtest/accounting_identity.pdf}.

\end{thebibliography}

\end{document}

